\chapter{Entropy, Perelman and the Uncertainty Principle}
\label{ch:entropy}

\section{Four entropies}

Clausius introduced entropy in 1865~\cite{clausius1865} as a state function whose differential is $\delta Q/T$ along reversible processes. Boltzmann gave it a microscopic meaning in 1872~\cite{boltzmann1872} through his $H$ functional $H(f) = \int f\log f$, and in 1877~\cite{boltzmann1877} by showing that the entropy of a macrostate is proportional to the logarithm of the number of microstates that realize it, a relation that Planck wrote as $S = k\log W$~\cite{planck1901}. Boltzmann's $H$ is the negative of the entropy: $H(f) = -h(f)$ for the differential entropy $h$ defined below, so it has the opposite sign to the Shannon entropy $H(p) = -\sum_jp_j\log p_j$ of Chapter~\ref{ch:axioms}. Gibbs in 1902~\cite{gibbs1902} defined the entropy of an ensemble as $-k\int\rho\log\rho$ on phase space. Shannon in 1948~\cite{shannon1948} defined the entropy of a source and proved that it is the minimal rate of lossless compression. For a density $u$ on $\R^n$ we write
\[
h(u) = -\int u\log u\dd x,\qquad J(u) = \int\frac{|\nabla u|^2}{u}\dd x = \int u\,|\nabla\log u|^2\dd x
\]
for the differential entropy and the Fisher information. Fisher~\cite{fisher1925} introduced in 1925 the information that a single observation carries about a parameter, and $J(u)$ is this information for the location family generated by $u$ (for $n>1$ the trace of the information matrix).

\section{De Bruijn's identity}

\begin{theorem}[de Bruijn]
\label{thm:debruijn}
Let $X$ have finite variance and let $Z\sim N(0,I_n)$ be independent. Let $u_t$ be the density of $X+\sqrt t Z$. Then for $t>0$,
\[
\frac{\dd}{\dd t}h(u_t) = \frac12J(u_t).
\]
\end{theorem}

\begin{proof}
For $t>0$ the density $u_t = \E\,g_t(\cdot-X)$, where $g_t$ is the $N(0,tI_n)$ density, is smooth and positive and solves $\partial_tu = \frac12\Delta u$. Jensen's inequality gives $u_t(x)\geq g_t(0)\exp(-\E|x-X|^2/2t)$, while $u_t\leq g_t(0)$, so $|\log u_t(x)|\leq C(1+|x|^2)$. The derivatives of $g_t$ are polynomials times $g_t$, so $u_t$, $\nabla u_t$ and $\Delta u_t$ are integrable against $1+|x|^2$, locally uniformly in $t$, and the Cauchy Schwarz inequality gives $J(u_t)\leq n/t$. This justifies the differentiation under the integral and the integrations by parts below. Then
\[
\begin{aligned}\frac{\dd}{\dd t}h(u_t) &= -\int(\partial_tu)(1+\log u)\dd x = -\frac12\int\Delta u\,\log u\dd x\\ &= \frac12\int\nabla u\cdot\nabla\log u\dd x = \frac12J(u_t),\end{aligned}
\]
where we used $\int\Delta u = 0$.
\end{proof}

The identity appears in Stam's paper~\cite{stam1959}, and an integrated form is due to Barron~\cite{barron1986}.

\begin{theorem}[Cram\'er Rao, Stam]
\label{thm:stam}
For a density $u$ on $\R$ with variance $\sigma^2$, $J(u)\geq1/\sigma^2$, with equality only for the Gaussian. For independent $X,Y$, Stam's inequality holds:
\[
\frac{1}{J(X+Y)}\geq\frac{1}{J(X)}+\frac{1}{J(Y)}.
\]
\end{theorem}

\begin{proof}
For the first claim, $\int(x-m)u'(x)\dd x = -1$ by integration by parts, and Cauchy Schwarz gives $1\leq\int(x-m)^2u\cdot\int(u')^2/u = \sigma^2J(u)$. For Stam's inequality~\cite{stam1959}, in the form of Blachman's proof~\cite{blachman1965}, the score of $X+Y$ is a conditional expectation of the scores of $X$ and of $Y$: $\rho_{X+Y}(X+Y) = \E[\rho_X(X)\mid X+Y] = \E[\rho_Y(Y)\mid X+Y]$, where $\rho = (\log u)'$. Hence for any $a+b=1$, $\rho_{X+Y} = \E[a\rho_X(X)+b\rho_Y(Y)\mid X+Y]$, and Jensen's inequality gives $J(X+Y)\leq\E\big[(a\rho_X(X)+b\rho_Y(Y))^2\big] = a^2J(X)+b^2J(Y)$, the cross term vanishing because $X$ and $Y$ are independent and the scores have mean zero. Optimizing over $a$ gives the claim.
\end{proof}

Integrating de Bruijn's identity along the heat flow and using Stam's inequality yields the entropy power inequality of Shannon and Stam~\cite{shannon1948,stam1959,blachman1965}, which for independent random vectors $X$ and $Y$ in $\R^n$ reads $e^{2h(X+Y)/n}\geq e^{2h(X)/n}+e^{2h(Y)/n}$, and the Gaussian logarithmic Sobolev inequality of Gross~\cite{gross1975}. We prove the second below and refer to~\cite{covertomas} for the first.

\begin{theorem}[Gross]
\label{thm:lsi}
Let $\gamma$ be the standard Gaussian measure on $\R^n$. For smooth $f$ with $\int f^2\dd\gamma = 1$,
\[
\int f^2\log f^2\dd\gamma\leq2\int|\nabla f|^2\dd\gamma.
\]
Equivalently, for every probability density $u$ on $\R^n$ with finite Fisher information,
\[
h(u)\geq\frac n2\log\Big(\frac{2\pi e\,n}{J(u)}\Big),
\]
with equality for isotropic Gaussians.
\end{theorem}

\begin{proof}[Proof of the Euclidean form]
Assume first that $u$ has finite second moment, and let $u_t$ be the density of $X+\sqrt tZ$ where $X\sim u$. Stam's inequality holds verbatim in $\R^n$ with $J$ the trace of the Fisher information matrix, since the proof of Theorem~\ref{thm:stam} only uses that the score of a sum is a conditional expectation. With $J(\sqrt tZ) = n/t$ it gives $J(u_t)^{-1}\geq J(u)^{-1}+t/n$, that is $J(u_t)\leq nJ(u)/(n+tJ(u))$~\cite{carlen1991}. By de Bruijn's identity in its integrated form~\cite{barron1986},
\[
\begin{aligned}h(u_T) - h(u) &= \frac12\int_0^TJ(u_t)\dd t\\ &\leq\frac12\int_0^T\frac{n\,J(u)}{n+tJ(u)}\dd t = \frac n2\log\Big(1+\frac{TJ(u)}{n}\Big).\end{aligned}
\]
Since conditioning reduces entropy, $h(u_T) = h(X+\sqrt TZ)\geq h(X+\sqrt TZ\mid X) = h(\sqrt TZ) = \frac n2\log(2\pi eT)$. Hence
\[
h(u)\geq\frac n2\log(2\pi eT)-\frac n2\log\Big(1+\frac{TJ(u)}{n}\Big) = \frac n2\log\frac{2\pi e\,nT}{n+TJ(u)},
\]
and letting $T\to\infty$ gives $h(u)\geq\frac n2\log(2\pi e\,n/J(u))$.

If $u$ has infinite second moment, let $\chi$ be smooth with $0\leq\chi\leq1$, $\chi = 1$ on the unit ball and $\chi = 0$ outside the ball of radius $2$, and put $\chi_R = \chi(\cdot/R)$, $m_R = \int\chi_R^2u\dd x$ and $u_R = \chi_R^2u/m_R$. Then $u_R$ has compact support, $m_R\to1$, and $J(u_R)^{1/2}\leq m_R^{-1/2}\big(J(u)^{1/2}+2\|\nabla\chi\|_\infty/R\big)$, because $J(u) = 4\int|\nabla\sqrt u|^2\dd x$ and $\int u\dd x = 1$. Finite Fisher information gives $\sqrt u\in H^1(\R^n)$, hence $u\in L^p$ for some $p>1$ and $h(u)>-\infty$. If $h(u)<\infty$, then $h(u_R)\to h(u)$ by dominated convergence, and the inequality for $u_R$ passes to the limit; if $h(u) = +\infty$ there is nothing to prove.
\end{proof}

\begin{proof}[Proof of the Gaussian form and of the equivalence]
Write $u = f^2\gamma$ and $q = \int|x|^2u\dd x$, and suppose that $q$ and $J(u)$ are finite (the general case follows by approximation). Since $\nabla\log u = 2\nabla f/f-x$ and $\int x\cdot\nabla(f^2)\dd\gamma = \int f^2(|x|^2-n)\dd\gamma = q-n$,
\[
\begin{aligned}J(u) &= 4\int|\nabla f|^2\dd\gamma+2n-q,\\ h(u) &= -\int f^2\log f^2\dd\gamma+\frac n2\log2\pi+\frac q2.\end{aligned}
\]
Adding,
\[
h(u)+\frac12J(u)-\frac n2\log2\pi-n = 2\int|\nabla f|^2\dd\gamma-\int f^2\log f^2\dd\gamma,
\]
and the Gaussian form says that the right side is nonnegative. Suppose the Euclidean form holds. Since $\log x\geq1-1/x$, we have $\frac n2\log\frac{n}{J(u)}\geq\frac n2-\frac12J(u)$, so $h(u)\geq\frac n2\log(2\pi e)+\frac n2\log\frac{n}{J(u)}\geq\frac n2\log2\pi+n-\frac12J(u)$, and the left side of the display is nonnegative. This is the Gaussian form. Conversely, suppose the Gaussian form holds. Then $h(v)+\frac12J(v)\geq\frac n2\log2\pi+n$ for every smooth positive density $v$ with finite second moment and finite Fisher information (the general case follows by approximation). Apply this to $v(x) = \lambda^nu(\lambda x)$, for which $h(v) = h(u)-n\log\lambda$ and $J(v) = \lambda^2J(u)$, and choose $\lambda^2 = n/J(u)$. This gives $h(u)+\frac n2\log\frac{J(u)}{n}+\frac n2\geq\frac n2\log2\pi+n$, which is the Euclidean form.
\end{proof}

The log Sobolev inequality strengthens the Poincar\'e inequality of Theorem~\ref{thm:gaussian_poincare}: it controls entropy rather than variance, and it implies exponential decay of relative entropy along the Ornstein Uhlenbeck flow, $\KL(\rho_t\|\gamma)\leq e^{-2t}\KL(\rho_0\|\gamma)$. Bakry and \'Emery~\cite{bakry1985} showed in 1985 that for the diffusion with generator $\Delta - \nabla V\cdot\nabla$ and invariant measure $\mu\propto e^{-V}$ the relative entropy decays exponentially whenever $\nabla^2V\geq K I$ with $K>0$, at the rate $\KL(\rho_t\|\mu)\leq e^{-2Kt}\KL(\rho_0\|\mu)$, and more generally on a Riemannian manifold whenever $\Ric+\nabla^2V\geq K$.

\section{The second derivative of entropy}

\begin{lemma}[McKean]
\label{lem:mckean}
Let $u_\tau$ solve $\partial_\tau u = \Delta u$ on $\R^n$ with $u_0$ a probability density of finite variance. Then for $\tau>0$,
\[
\frac{\dd}{\dd\tau}h(u_\tau) = J(u_\tau),\qquad\frac{\dd}{\dd\tau}J(u_\tau) = -2\int u_\tau\,\big\|\nabla^2\log u_\tau\big\|^2\dd x,
\]
where $\|\cdot\|$ is the Hilbert Schmidt norm.
\end{lemma}

\begin{proof}
The first identity is Theorem~\ref{thm:debruijn} with time rescaled by $\tau = t/2$. For the second write $\ell = \log u$, so that $\partial_\tau\ell = \Delta\ell+|\nabla\ell|^2$. Then
\[
\frac{\dd}{\dd\tau}\int u|\nabla\ell|^2 = \int\Delta u\,|\nabla\ell|^2 + 2\int u\,\nabla\ell\cdot\nabla(\Delta\ell+|\nabla\ell|^2).
\]
On flat space Bochner's identity reads $\frac12\Delta|\nabla\ell|^2 = \|\nabla^2\ell\|^2+\nabla\ell\cdot\nabla\Delta\ell$. Integrating by parts twice, the first term is
\[
\int\Delta u\,|\nabla\ell|^2 = \int u\,\Delta|\nabla\ell|^2 = 2\int u\big(\|\nabla^2\ell\|^2+\nabla\ell\cdot\nabla\Delta\ell\big).
\]
The second term splits as $2\int u\nabla\ell\cdot\nabla\Delta\ell$ plus
\[
\begin{aligned}2\int u\,\nabla\ell\cdot\nabla|\nabla\ell|^2 &= 2\int\nabla u\cdot\nabla|\nabla\ell|^2 = -2\int u\,\Delta|\nabla\ell|^2\\ &= -4\int u\big(\|\nabla^2\ell\|^2+\nabla\ell\cdot\nabla\Delta\ell\big),\end{aligned}
\]
using $u\nabla\ell = \nabla u$. Adding the three contributions, the terms in $\nabla\ell\cdot\nabla\Delta\ell$ cancel ($2+2-4 = 0$) and the terms in $\|\nabla^2\ell\|^2$ give $2-4 = -2$, which is the claim. The decay of $u$ and its derivatives for $\tau>0$ justifies the integrations by parts.
\end{proof}

McKean~\cite{mckean1966} found the second identity in one dimension, and Bochner's identity gives it in $\R^n$. On a Riemannian manifold Bochner's identity acquires the term $\Ric(\nabla\ell,\nabla\ell)$, and this term is where curvature enters the entropy computations of geometric analysis.

\section{Perelman's entropy}

Perelman's functional, recalled in Chapter~\ref{ch:poincare}, is
\begin{gather*}
\mathcal{W}(g,f,\tau) = \int_M\big[\tau(R+|\nabla f|^2)+f-n\big]\,u\dd\,\mathrm{vol}, \\
u = (4\pi\tau)^{-n/2}e^{-f},\quad\int_Mu\dd\,\mathrm{vol} = 1.
\end{gather*}
Under the coupled system $\partial_tg = -2\Ric$, $\partial_t\tau = -1$, $\partial_tf = -\Delta f+|\nabla f|^2-R+\frac{n}{2\tau}$, Perelman proved~\cite{perelman2002}
\[
\frac{\dd}{\dd t}\mathcal{W} = 2\tau\int_M\Big\|\Ric+\nabla^2f-\frac{g}{2\tau}\Big\|^2u\dd\,\mathrm{vol}\geq0,
\]
with equality exactly on shrinking gradient solitons. The function $u$ solves the conjugate heat equation, which runs backward in $t$, so forward in $\tau$.

On flat $\R^n$ the metric does not move, $R=0$, and the formula becomes a statement about the heat equation. The proof below is analytic, by Bochner's identity; Section~\ref{sec:scalemixtures} describes a probabilistic proof for scale mixtures.

\begin{theorem}[Perelman's entropy on flat space]
\label{thm:perelman_flat}
Let $u_\tau$ solve $\partial_\tau u = \Delta u$ on $\R^n$ with $u_0$ a probability density of finite variance, and define
\[
W(\tau) = \tau J(u_\tau)+h(u_\tau)-\frac n2\log(4\pi\tau)-n.
\]
Then $W$ is nonincreasing on $(0,\infty)$, with
\[
\frac{\dd W}{\dd\tau} = -2\tau\int u_\tau\Big\|\nabla^2\log u_\tau+\frac{I}{2\tau}\Big\|^2\dd x\leq0,
\]
and $W(\tau)\to0$ as $\tau\to\infty$. Hence $W(\tau)\geq0$ for every $\tau$, with equality if and only if $u_\tau$ is a Gaussian density with covariance $2\tau I$. Moreover, the inequality $\tau J(v)+h(v)-\frac n2\log(4\pi\tau)-n\geq0$, for all probability densities $v$ and all $\tau>0$, is equivalent to the Euclidean logarithmic Sobolev inequality.
\end{theorem}

\begin{proof}
With $f = -\log u-\frac n2\log(4\pi\tau)$ we have $\int fu = h(u)-\frac n2\log(4\pi\tau)$ and $\int|\nabla f|^2u = J(u)$, so $W$ is Perelman's $\mathcal{W}$ with $R = 0$. Write $\ell = \log u_\tau$. By Lemma~\ref{lem:mckean},
\[
W' = J+\tau J'+J-\frac{n}{2\tau} = 2J-2\tau\int u\|\nabla^2\ell\|^2-\frac{n}{2\tau}.
\]
Expanding the square and using $\int u\,\Delta\ell = -\int\nabla u\cdot\nabla\ell = -J$,
\[
\begin{aligned}-2\tau\int u\Big\|\nabla^2\ell+\frac{I}{2\tau}\Big\|^2 &= -2\tau\int u\|\nabla^2\ell\|^2-2\int u\,\Delta\ell-\frac{n}{2\tau}\\ &= -2\tau\int u\|\nabla^2\ell\|^2+2J-\frac{n}{2\tau},\end{aligned}
\]
which is $W'$. For the limit, let $X\sim u_0$ and $\sigma^2 = \E|X-\E X|^2$, so that $u_\tau$ is the density of $X+\sqrt{2\tau}Z$ and has variance $\sigma^2+2n\tau$. Stam's inequality gives $J(u_\tau)\leq n/(2\tau)$, since $J(\sqrt{2\tau}Z) = n/(2\tau)$, and the Cram\'er Rao inequality in $\R^n$, $J(u)\geq n^2/\E|X-\E X|^2$, which follows as in the proof of Theorem~\ref{thm:stam}, gives $J(u_\tau)\geq n^2/(\sigma^2+2n\tau)$. Hence $\tau J(u_\tau)\to n/2$. Since conditioning reduces entropy, $h(u_\tau)\geq h(\sqrt{2\tau}Z) = \frac n2\log(4\pi e\tau)$, and since the Gaussian maximizes entropy for given variance, $h(u_\tau)\leq\frac n2\log\big(2\pi e(\sigma^2+2n\tau)/n\big)$. Therefore
\[
\frac n2\leq h(u_\tau)-\frac n2\log(4\pi\tau)\leq\frac n2+\frac n2\log\Big(1+\frac{\sigma^2}{2n\tau}\Big),
\]
and the right side tends to $\frac n2$, so $W(\tau)\to\frac n2+\frac n2-n = 0$. Since $W$ is nonincreasing, $W(\tau)\geq0$ for every $\tau$.

Suppose $W(\tau_0) = 0$. Then $W = 0$ and $W' = 0$ on $(\tau_0,\infty)$, so $\nabla^2\log u_s = -I/(2s)$ for $s>\tau_0$, and $u_s$ is a Gaussian density with covariance $2sI$. Dividing the characteristic function of $u_s$ by that of $N(0,2(s-\tau_0)I)$ shows that $u_{\tau_0}$ is a Gaussian density with covariance $2\tau_0I$. Conversely, for the Gaussian with covariance $2\tau I$ we have $J = n/(2\tau)$ and $h = \frac n2\log(4\pi e\tau)$, so $W = 0$. For a probability density $u_0$ the equality cannot occur, because $u_\tau$ then has covariance strictly larger than $2\tau I$; it is attained by the fundamental solution, whose initial datum is a point mass.

Finally, fix a probability density $v$. The function $\tau\mapsto\tau J(v)-\frac n2\log(4\pi\tau)$ is minimized at $\tau = n/(2J(v))$, and the minimum over $\tau$ of $\tau J(v)+h(v)-\frac n2\log(4\pi\tau)-n$ is
\[
h(v)-\frac n2\log\Big(\frac{2\pi e\,n}{J(v)}\Big).
\]
Hence the inequality of the theorem holds for all $v$ and $\tau$ exactly when this expression is nonnegative for all $v$, which is the Euclidean logarithmic Sobolev inequality of Theorem~\ref{thm:lsi}.
\end{proof}

Perelman's entropy is a deficit rather than a distance: a scale invariant combination of entropy and Fisher information that is nonnegative along every solution of the heat equation with finite variance and vanishes only on the fundamental solution. For the linear heat equation on a closed manifold the same functional is due to Ni~\cite{ni2004}, who showed that it is nonincreasing when $\Ric\geq0$.

\subsection{Scale mixtures and subordinated Brownian motion}
\label{sec:scalemixtures}

The flat theorem has a transparent probabilistic structure when the initial law is a scale mixture. Suppose $u_0(x) = \int\varphi_v(x)\,G(\dd v)$, where $\varphi_v$ is the centered Gaussian density with covariance $vI$. Then $u_\tau = \int\varphi_{v+2\tau}\,G(\dd v)$: the heat flow acts on the mixing law by translation. In terms of processes, if $X_0 = B_V$ is Brownian motion subordinated by an independent random time $V\sim G$, then $X_\tau\eqd B_{V+2\tau}$. Entropy, Fisher information and therefore $W$ become functionals of the translated mixing law, and the convergence $W(\tau)\to0$ expresses the fact that the relative spread of $V+2\tau$ vanishes as $\tau\to\infty$. In our paper on Perelman's entropy~\cite{polsonsokolovperelman} we use this representation, together with the Yor identities of Chapter~\ref{ch:brownian} for the mixing variables, to give a probabilistic proof of the monotonicity for the class of scale mixtures and to interpret the soliton condition as a statement about the mixing law. The geometric content of Perelman's formula, the Ricci term, is then the correction required when the Brownian motion runs on a curved space.

\section{The uncertainty principle as an information inequality}

Let $\psi\in L^2(\R)$ be a normalized wave function with position density $u = |\psi|^2$. The momentum operator is $P = -i\hbar\,\dd/\dd x$, and $\langle P^2\rangle = \hbar^2\int|\psi'|^2\dd x$. Write $\psi = \sqrt u\,e^{iS/\hbar}$. Then
\[
\int|\psi'|^2 = \int\Big(\big(\sqrt u\big)'\Big)^2+\frac{1}{\hbar^2}\int u\,(S')^2 = \frac14J(u)+\frac{1}{\hbar^2}\int u(S')^2.
\]
Hence the momentum variance satisfies
\[
\Var(P)\geq\frac{\hbar^2}{4}J(u)-\Big(\int u\,S'\Big)^2+\int u(S')^2\geq\frac{\hbar^2}{4}J(u),
\]
by Cauchy Schwarz. Combining with the Cram\'er Rao bound $J(u)\geq1/\Var(X)$ of Theorem~\ref{thm:stam} gives
\[
\Var(X)\,\Var(P)\geq\frac{\hbar^2}{4}.
\]
Equality holds exactly for Gaussian wave packets with linear phase. The same computation in $n$ dimensions gives $\langle|P|^2\rangle\geq\frac{\hbar^2}{4}J(u)$, which is the starting point of Nelson's stochastic mechanics in Chapter~\ref{ch:quantum}.

A stronger statement replaces variances by entropies. Hirschman~\cite{hirschman1957} conjectured, and Beckner~\cite{beckner1975} and Bia\l{}ynicki Birula and Mycielski~\cite{bbm1975} proved, that for $\psi$ normalized in $L^2(\R^n)$ with Fourier transform $\hat\psi(\xi) = (2\pi)^{-n/2}\int e^{-ix\cdot\xi}\psi(x)\dd x$,
\[
h\big(|\psi|^2\big)+h\big(|\hat\psi|^2\big)\geq n\log(\pi e).
\]
The entropic uncertainty principle implies Heisenberg's because entropy is bounded by the logarithm of variance, and it is sharp for Gaussians. With the convention $e^{-2\pi ix\cdot\xi}$ for the Fourier transform the sharp constant is $n\log(e/2)$ instead. The proof of the inequality uses the sharp Hausdorff Young inequality of Beckner~\cite{beckner1975}.

\section{Entropy as a gradient flow}

In 1998 Jordan, Kinderlehrer and Otto~\cite{jko1998} discovered that the Fokker Planck equation
\begin{equation}
\partial_t\rho = \nabla\cdot(\rho\nabla V)+\Delta\rho
\label{eq:fp}
\end{equation}
is the gradient flow of the free energy $\mathcal{F}(\rho) = \int V\rho+\int\rho\log\rho$ with respect to the Wasserstein distance
\[
W_2(\mu,\nu)^2 = \inf_{\pi\in\Pi(\mu,\nu)}\int|x-y|^2\,\pi(\dd x,\dd y).
\]
Precisely, for a time step $\eta>0$ define recursively
\[
\rho_{k+1}^\eta = \argmin_\rho\Big\{\frac{1}{2\eta}W_2(\rho,\rho_k^\eta)^2+\mathcal{F}(\rho)\Big\}.
\]

\begin{theorem}[Jordan, Kinderlehrer and Otto]
\label{thm:jko}
Let $V\geq0$ be smooth with $|\nabla V|\leq C(V+1)$, and let $\rho_0$ have finite second moment and finite free energy. The piecewise constant interpolation of $(\rho^\eta_k)$ converges, as $\eta\to0$, weakly in $L^1$ for every $t>0$, to the unique solution of \eqref{eq:fp} with initial datum $\rho_0$.
\end{theorem}

We quote the theorem without proof; see Jordan, Kinderlehrer and Otto~\cite[Theorem~5.1]{jko1998}. The scheme is a minimizing movement. The heat equation is the case $V = 0$, the steepest descent of $\int\rho\log\rho$, that is the steepest ascent of Boltzmann's entropy, in the geometry of optimal transport. Otto's formal Riemannian calculus on the space of probability measures~\cite{otto2001} makes this precise: the tangent space at $\rho$ consists of velocity fields $\nabla\phi$ with norm $\int\rho|\nabla\phi|^2$, and the gradient of $\mathcal{F}$ is $-\nabla\cdot(\rho\nabla\frac{\delta\mathcal{F}}{\delta\rho})$.

The implications for this book are threefold. First, the dissipation of free energy along \eqref{eq:fp} equals the relative Fisher information $\int\rho|\nabla\log(\rho/e^{-V})|^2$, which is de Bruijn's identity in transport language. Second, the uniform convexity of $\mathcal{F}$ along Wasserstein geodesics, McCann's displacement convexity~\cite{mccann1997}, implies exponential convergence and the log Sobolev and Talagrand inequalities, unified in the HWI inequality of Otto and Villani~\cite{ottovillani2000}. Third, von Renesse and Sturm~\cite{vonrenesse2005} showed that a smooth Riemannian manifold has $\Ric\geq K$ if and only if Boltzmann's functional $H(\rho) = \int\rho\log\rho\,\dd\,\mathrm{vol}$ is $K$ convex along Wasserstein geodesics of probability measures on it, one direction having been proved earlier by Cordero Erausquin, McCann and Schmuckenschl\"ager~\cite{corderoerausquin2001}, and Lott and Villani~\cite{lottvillani2009} and Sturm~\cite{sturm2006} used this property to define lower bounds on the Ricci curvature of metric measure spaces. Ricci curvature, the quantity that drives Perelman's flow, is a property of entropy and transport.

For inference, Theorem~\ref{thm:jko} says that Langevin sampling, whose law evolves by \eqref{eq:fp} with $V = -\log\pi$, is gradient descent of the Kullback Leibler divergence to the posterior $\pi$ in the Wasserstein geometry, provided $V$ is smooth and bounded below and, after adding a constant so that $V\geq0$, satisfies the growth condition $|\nabla V|\leq C(V+1)$ of the theorem, and the initial law has finite second moment and finite free energy. Variational inference, Stein variational gradient descent and diffusion based generative models are all variations on this flow, and Chapter~\ref{ch:gbc} returns to them.

\section{Concentration from the log Sobolev inequality}

The log Sobolev inequality implies Gaussian concentration by a short argument due to Herbst (see Ledoux~\cite{ledoux2001}).

\begin{theorem}[Herbst]
Let $\mu$ satisfy $\Ent_\mu(f^2)\leq2C\int|\nabla f|^2\dd\mu$ for all smooth $f$, where $\Ent_\mu(g)=\int g\log g\dd\mu-\int g\dd\mu\,\log\int g\dd\mu$ for $g\ge0$. Then every $1$ Lipschitz function $F$ satisfies
\[
\mu\Big(F-\int F\dd\mu\geq r\Big)\leq e^{-r^2/2C},\qquad r\geq0.
\]
\end{theorem}

\begin{proof}
Let $M(\lambda) = \int e^{\lambda F}\dd\mu$ and apply the inequality to $f = e^{\lambda F/2}$. Since $|\nabla f|^2\leq\frac{\lambda^2}{4}e^{\lambda F}$, we get $\lambda M'(\lambda)-M(\lambda)\log M(\lambda)\leq\frac{C\lambda^2}{2}M(\lambda)$. Dividing by $\lambda^2M$ shows that $K(\lambda) = \lambda^{-1}\log M(\lambda)$ satisfies $K'(\lambda)\leq C/2$, and $K(0^+) = \int F\dd\mu$. Hence $\log M(\lambda)\leq\lambda\int F+C\lambda^2/2$, and Chebyshev's exponential inequality with $\lambda = r/C$ gives the bound.
\end{proof}

For the standard Gaussian measure $C = 1$, and the theorem says that a Lipschitz function of many independent Gaussian variables is as concentrated as a single Gaussian, with no dependence on dimension. This is one mechanism behind the determinism of the macroscopic observables of a gas with many particles: if the Gibbs measure satisfies a log Sobolev inequality with a constant that does not grow with the number of particles, then a macroscopic observable, a Lipschitz function of the microstate with small Lipschitz constant, barely fluctuates. Whether a given Gibbs measure satisfies such an inequality is a question about the interaction.

\section{The I MMSE identity}\label{sec:immse}

De Bruijn's identity has an estimation theoretic twin found by Guo, Shamai and Verd\'u in 2005~\cite{guo2005}. Let $Y = \sqrt{s}\,X+Z$ with $Z$ standard Gaussian independent of $X$ and $s>0$ the signal to noise ratio. Let $\mathrm{mmse}(s) = \E|X-\E[X\mid Y]|^2$.

\begin{theorem}[Guo, Shamai and Verd\'u]
\label{thm:immse}
For $X$ with finite second moment,
\[
\frac{\dd}{\dd s}I(X;Y) = \frac12\,\mathrm{mmse}(s),
\]
where the mutual information is measured in nats.
\end{theorem}

\begin{proof}[Sketch]
Write $I(X;Y) = h(Y)-h(Z)$ and rescale: $Y/\sqrt s = X+Z/\sqrt s$, so $h(Y) = h(X+\sqrt tZ')+\frac n2\log s$ with $t = 1/s$ and $Z'$ standard. De Bruijn's identity gives $\frac{\dd}{\dd t}h(X+\sqrt tZ') = \frac12J(X+\sqrt tZ')$, and Tweedie's formula with the Fisher information identity $J(X+\sqrt tZ') = \frac nt-\frac{1}{t^2}\mathrm{mmse}$ in the rescaled problem converts the derivative in $t$ into the stated derivative in $s$.
\end{proof}

In the inference problems of Chapter~\ref{ch:spins}, where the replica formula gives the mutual information, Theorem~\ref{thm:immse} turns the derivative of the mutual information in the signal to noise ratio into the Bayes optimal error, and so locates its phase transitions.

\section{Transport inequalities}

Talagrand proved in 1996~\cite{talagrand1996} that the standard Gaussian measure satisfies the transport entropy inequality
\[
W_2(\nu,\gamma)^2\leq2\,\KL(\nu\,\|\,\gamma)
\]
for every probability measure $\nu$. Otto and Villani~\cite{ottovillani2000} showed that if a probability measure $\mu$ satisfies the log Sobolev inequality with constant $C$, then it satisfies Talagrand's inequality $W_2(\nu,\mu)^2\leq2C\,\KL(\nu\,\|\,\mu)$ with the same constant, and they introduced the HWI inequality, which for the standard Gaussian measure reads
\[
\KL(\nu\|\gamma)\leq W_2(\nu,\gamma)\sqrt{J(\nu\|\gamma)}-\frac12W_2(\nu,\gamma)^2,
\]
where $J(\nu\|\gamma) = \int|\nabla\log\frac{\dd\nu}{\dd\gamma}|^2\dd\nu$ is the relative Fisher information. The three letters of HWI stand for relative entropy ($\KL$), transport distance ($W_2$) and Fisher information ($J$), the three quantities of this chapter. For Bayesian computation, suppose that the posterior $\pi$ satisfies the log Sobolev inequality with constant $C$. Then Talagrand's inequality bounds the Wasserstein error of an approximate posterior $\nu$ by its Kullback Leibler error, $W_2(\nu,\pi)^2\leq2C\,\KL(\nu\,\|\,\pi)$, and the Kullback Leibler error of the Langevin sampler at time $t$ is at most $e^{-2t/C}$ times its initial value.

\section{Boltzmann, Shannon and the second law of inference}

Boltzmann's $H$ and Shannon's entropy differ by a sign and a constant. Let $\pi$ be a prior and $\pi_y$ a posterior. Averaged over the predictive distribution of $y$,
\[
\E_y\big[\KL(\pi_y\,\|\,\pi)\big] = I(\theta;Y)\geq0,
\]
so the expected information gain is the mutual information, and the expected posterior entropy is never larger than the prior entropy. This is the second law of inference, with the arrow reversed: observation decreases expected entropy about the parameter, at the price, by Landauer's principle of Chapter~\ref{ch:kolmogorov}, of increasing the entropy of the environment when the observation is recorded and later erased. The parallel with Boltzmann's $H$ theorem is only partial. For the Boltzmann equation $H$ decreases along the solution, whereas the posterior entropy can increase for particular data, when a surprising observation spreads the posterior out; only its average is monotone.

\section*{Exercises}
\begin{enumerate}
\item For an equal mixture of $N(\pm m,1)$, compute $h$, $J$ and $W(\tau)$ numerically for several $m$ and check that $W\geq0$.
\item Show that Stam's inequality is an equality for independent Gaussian $X$ and $Y$.
\item Show that the log Sobolev inequality implies the Poincar\'e inequality by applying it to $f = 1+\varepsilon g$ and expanding in $\varepsilon$.
\item Derive $\Var(X)\Var(P)\geq\hbar^2/4$ from the entropic uncertainty principle and the maximum entropy property of Gaussians.
\item Show that the JKO scheme with $V = 0$ and Gaussian initial data produces Gaussian iterates, and compute their variances.
\end{enumerate}


\section*{Notes}

Cover and Thomas~\cite{covertomas} treat de Bruijn's identity, Stam's inequality and the entropy power inequality. Carlen~\cite{carlen1991} derived the log Sobolev inequality from Stam's inequality. Villani's \emph{Topics in Optimal Transportation}~\cite{villani2003} and \emph{Optimal Transport, Old and New}~\cite{villani2009} contain the JKO theorem, Otto calculus and the Lott Sturm Villani theory. Perelman's entropy is explained in detail by Kleiner and Lott~\cite{kleinerlott2008}, and Morgan and Tian~\cite{morgan2007} write out the proof of the Poincar\'e conjecture.
